\chapter{The Wu Fortran Solver}
\label{chap:wu_solver}

This chapter describes the numerical core: W.~Wu's finite-difference
piecewise PV-inversion code, driven from Python (\texttt{wu/steps/})
while the solver itself remains the original Fortran-77.  The three
programs live in \texttt{wu/fortran/} and are compiled with
\texttt{gfortran -std=legacy -O2 -fno-automatic}.

\section{Grid and Vertical Levels}
\label{sec:wu_grid}

The solver runs on an extended Northern-Hemisphere window
(\texttt{config.py}, \texttt{wu/wu\_config.yaml}):

\begin{center}\small
\begin{tabular}{lll}
  \toprule
  Quantity & Symbol / code & Value \\
  \midrule
  Longitudes & $N_X$ & $134$ \\
  Latitudes  & $N_Y$ & $51$ \\
  Levels     & $N_W=N_L$ & $9$ \\
  Latitude row & $\phi_i = 85.5 - 1.5\,i$ & $85.5^\circ\!\to\!10.5^\circ$ (N$\to$S) \\
  Longitude col & $\lambda_j = 120 + 1.5\,j$ & $120^\circ\mathrm{E}\!\to\!319.5^\circ\mathrm{E}$ ($=40.5^\circ$W) \\
  Pressure levels & PLEVS [hPa] & $1000,850,700,500,400,$ \\
                  &             & $300,250,200,100$ \\
  Earth radius & $a=2\times10^{7}/\pi$ & $6.366\times10^{6}$ m \\
  Grid spacing & $\Delta p=\Delta l = a\,\mathrm{rad}(1.5^\circ)$ & $1.667\times10^{5}$ m \\
  \bottomrule
\end{tabular}
\end{center}

\begin{warnbox}[Fortran \texttt{PARAMETER} blocks are hard-wired]
$N_X,N_Y,N_W$ appear as compile-time \texttt{PARAMETER}s inside
\emph{every} \texttt{.f} file (and inside each subroutine).  Changing the
grid requires editing and recompiling the Fortran --- not just
\texttt{config.py}.  The current 9-level build sets \texttt{NW=9},
\texttt{NL=9}, and a 9-value \texttt{DATA PR} array in all three sources.
The Python drivers pass solver \emph{parameters} (relaxation factors,
piece definitions) over \texttt{stdin}; they cannot change array
dimensions.
\end{warnbox}

\section{The Three Programs and Four Passes}
\label{sec:wu_passes}

\begin{center}\small
\begin{tabular}{llll}
  \toprule
  Pass & Program & Subroutine & Produces \\
  \midrule
  A & \texttt{pvpialln}  & ---    & \texttt{meanq}, \texttt{meanh} (mean $Q$, $\theta_b$, $\psi$) \\
  B & \texttt{pvpialln}  & ---    & \texttt{event\_q.out}, \texttt{event\_h.out} \\
  C & \texttt{qinvert21} & BALNC  & \texttt{event\_bal.out} (total balanced $\psi$, $\Phi$) \\
  D & \texttt{qinvertp21}& BALP   & \texttt{event\_pert.out} (3-piece $\psi'$) \\
  \bottomrule
\end{tabular}
\end{center}

\paragraph{Pass A/B --- \texttt{pvpialln}.}  Computes the boundary
potential temperatures $\theta_b^{\text{bot}},\theta_b^{\text{top}}$
(\Cref{sec:boundary_theta}), the Ertel-like PV on the interior levels
($K=2\ldots N_W{-}1$, i.e.\ 850--200\,hPa), the balanced streamfunction
$\psi$ from relative-vorticity inversion, and the geopotential.  Pass~A
runs on the 30-year climatological mean state; Pass~B on the
2025-01-08 event state.  The \texttt{q.out} layout is
$[\theta_b^{\text{bot}},\theta_b^{\text{top}},Q(K{=}2..N_W{-}1)]$.

\paragraph{Pass C --- \texttt{qinvert21} (BALNC).}  The full
\emph{total} balanced inversion: a 2-D Helmholtz problem for $\psi$ on
each level coupled to a 3-D PV--$\psi$ relation under gradient-wind
balance, by SOR with under-relaxation between $\psi$ and $\Phi$.  The
boundary $\theta$ enters as the inhomogeneous Neumann condition of
\Cref{sec:boundary_theta}.

\paragraph{Pass D --- \texttt{qinvertp21} (BALP).}  Partitions the PV
\emph{perturbation} $\dq=q_{\text{event}}-q_{\text{mean}}$ into three
pieces and inverts each independently to its perturbation streamfunction
$\psi'$.  In the BALP convention a piece is a list \texttt{QLV} of level
indices, where the value \texttt{1} activates the \emph{lower-boundary}
$\theta$, \texttt{NL} the \emph{upper-boundary} $\theta$, and
$2..N_L{-}1$ the interior PV at that level:

\begin{center}\small
\begin{tabular}{llll}
  \toprule
  Piece & \texttt{QLV} & Activates & Pressures \\
  \midrule
  \textbf{surface} & $\{1\}$              & 1000\,hPa boundary $\theta$ only & --- \\
  \textbf{lower}   & $\{2,3\}$            & interior PV & $850,700$ hPa \\
  \textbf{upper}   & $\{4,5,6,7,8,9\}$    & interior PV $+$ top $\theta$ ($\texttt{9}{=}N_L$) & $500$--$200$ hPa $+\,100\,\theta$ \\
  \bottomrule
\end{tabular}
\end{center}

The three \texttt{QLV} lists cover every field
($\theta_b^{\text{bot}}$, interior PV at $K{=}2..8$, $\theta_b^{\text{top}}$)
exactly once, which is the precondition for the additivity
\cref{eq:additivity}.  The stdin actually sent is
\texttt{1,1}\,/\,\texttt{2,2,3}\,/\,\texttt{5,4,5,6,7,8,9}.

\begin{keybox}[The surface piece is the surface heating]
Because \texttt{QLV}=$\{1\}$ activates only the lower-boundary $\theta$,
the \textbf{surface piece is the inversion of the surface
warm/cold-$\theta$ anomaly} (the Bretherton PV sheet,
\Cref{sec:boundary_theta}) --- not of any interior 1000\,hPa PV (there is
none).  This isolates the surface thermal forcing from the
free-tropospheric PV, mirroring the reference inv3d partition.
\end{keybox}

\section{Solver Parameters}
\label{sec:wu_params}

Parameters are passed over \texttt{stdin} from \texttt{wu\_config.yaml}:

\begin{center}\small
\begin{tabular}{llll}
  \toprule
  Param & Pass C & Pass D & Role \\
  \midrule
  \texttt{OMEGS} & $1.4$  & $1.85$ & SOR factor, $\psi$ sweep \\
  \texttt{OMEGH} & $1.4$  & $1.4$  & SOR factor, $\Phi/H$ sweep \\
  \texttt{PART}  & $0.05$ & $0.05$ & under-relaxation $\psi\!\leftrightarrow\!\Phi$ \\
  \texttt{THRSH} & $0.01$ & $0.1$  & convergence threshold \\
  \texttt{INLIN} & ---    & $1$    & \textbf{nonlinear} balance (additivity) \\
  \texttt{IBC}   & ---    & $0$    & homogeneous Dirichlet lateral walls \\
  \bottomrule
\end{tabular}
\end{center}

\noindent \texttt{qinvertp21\_94.f} \emph{hardcodes} its iteration caps
and PV floor ($\texttt{MAX}=8000$, $\texttt{MAXT}=2000$,
$\texttt{QMIN}=10^{-5}$); the YAML does not set them.  The
non-dimensional Exner increment is $\texttt{DPI}=500/N_L=55.56$
(auto-scaled with $N_L$).

\section{Convergence and Closure (current behaviour)}
\label{sec:wu_converge}

\begin{itemize}\setlength\itemsep{2pt}
  \item \textbf{Pass D} converges cleanly: all three pieces report
        \texttt{TOTAL CONVERGENCE}, no NaNs, with physical magnitudes
        (the upper piece dominates the 250\,hPa response).
  \item \textbf{Pass C} prints \texttt{PHI CONVERGED} (per-cycle residual
        $\sim 10^{-4}$, below threshold) and then a \emph{benign}
        ``started diverging'' from \texttt{qinvert21}'s outer-coupling
        heuristic (it fires when an outer cycle needs $>\!10$ more inner
        iterations after $>\!30$ cycles).  The written
        \texttt{event\_bal.out} is the per-cycle-converged field
        (identical at $\omega=1.2$ and $1.4$); it is used as the total
        reference for Pass~D.
  \item \textbf{Top boundary 100\,hPa is non-uniform} ($\Delta\sigma$
        jumps $0.05\!\to\!0.10$ from 200 to 100\,hPa).  This is accepted:
        the 9-level grid closes the piecewise sum substantially better
        (linear-closure $R^2\!\approx\!0.44$ at 250\,hPa) than a
        uniform-top 8-level grid ($R^2\!\approx\!-0.2$), and all
        analysis of interest is at $\le 200$\,hPa.
\end{itemize}
